% Section 6.2, from lectures/L06/appendix/b_debugging.md.

\section{Debugging: Fatal Errors, Sanitizers, GDB and Coverage}
\label{c6:sec:debuggingfatalerro}\label{c6:app:b}
This section covers what happens when a Zephyr program goes wrong, and the tools that find out why.
Some of them are Zephyr’s own and work the same on a board: fatal errors, the handler that decides
what one costs, and assertions. The others are the host’s, and are much of the reason to keep code
runnable on \code{native_sim}: AddressSanitizer, UndefinedBehaviorSanitizer, GDB and gcov. By the
end you should be able to explain every line the \file{lectures/L06/examples/debugging} example
prints, in each of its configurations.

\subsection{Fatal errors}

\subsubsection{What makes an error fatal}
A \textbf{fatal error} is one the running code cannot continue from. The kernel names the reason
with a number:

\begin{booktable}{@{}llL@{}}
  \thead{Reason} & \thead{Value} & \thead{Raised by} \\
  \midrule
  \code{K_ERR_CPU_EXCEPTION} & 0 & The CPU: a bus fault, an illegal instruction, a division by zero,
    on a board. \\
  \code{K_ERR_SPURIOUS_IRQ} & 1 & An interrupt that has no handler. \\
  \code{K_ERR_STACK_CHK_FAIL} & 2 & A stack overflow check. \\
  \code{K_ERR_KERNEL_OOPS} & 3 & \code{k_oops()}: \emph{this thread} cannot go on. \\
  \code{K_ERR_KERNEL_PANIC} & 4 & \code{k_panic()}, and every failed assertion: \emph{the system}
    cannot go on. \\
\end{booktable}

On \code{native_sim} there is no CPU exception to raise, since the code runs natively; a genuine
segmentation fault kills the host process like any other Linux program’s. Oopses, panics and
assertions work exactly as on a board.

\textbf{\code{k_oops()} and \code{k_panic()} are for code that has found something impossible}, not
for ordinary errors. A sensor that does not answer is an error to return; a reading below absolute
zero, from a sensor the code has already validated, is a broken assumption, and an oops is how a
thread gives up on it.

\subsubsection{What the kernel does}
Every fatal error goes through the same path. The kernel logs what happened, and then calls
\textbf{\code{k_sys_fatal_error_handler()}} with the reason:

\begin{console}
[00:00:00.000,000] <err> app: Impossible reading: -300 C
@ WEST_TOPDIR/lectures/L06/exercises/exercise3/src/main.c:29
[00:00:00.000,000] <err> os: >>> ZEPHYR FATAL ERROR 3: Kernel oops on CPU 0
[00:00:00.000,000] <err> os: Current thread: 0x40f6c0 (sampler)
\end{console}

\begin{itemize}
  \item \textbf{\code{@ file:line}} is printed on architectures that have no CPU exception to raise for an
    oops, \code{native_sim} among them. On a board, the exception’s registers say where instead.
  \item \textbf{\code{ZEPHYR FATAL ERROR 3: Kernel oops}} gives the reason’s number and name. Twister looks
    for this exact string to tell a crashed test from a failed one.
  \item \textbf{\code{Current thread}} names the thread, if \code{CONFIG_THREAD_NAME=y}; otherwise only its
    address.
\end{itemize}

All of it is logged, so \textbf{without \code{CONFIG_LOG=y} a fatal error prints nothing} but what
the handler prints itself.

\subsubsection{The handler}
The kernel’s own \code{k_sys_fatal_error_handler()} is a weak function that halts the system,
whatever the reason:

\begin{console}
[00:00:00.000,000] <err> os: Halting system
Exiting due to fatal error
\end{console}

\noindent On a board, halting means an endless loop with interrupts off, until the watchdog or a
person resets the device. On \code{native_sim}, it means that the process exits with \textbf{status
1}, after printing \code{Exiting due to fatal error} on stderr.

An application that wants a different policy defines the function itself, which replaces the weak
one:

\begin{ccode}
#include <zephyr/fatal.h>
#include <zephyr/logging/log_ctrl.h>

void k_sys_fatal_error_handler(unsigned int reason, const struct arch_esf *esf)
{
	ARG_UNUSED(esf);

	LOG_PANIC();
	LOG_ERR("Fatal error %u in thread '%s'", reason, k_thread_name_get(k_current_get()));

	if (reason == K_ERR_KERNEL_OOPS) {
		return;
	}
	k_fatal_halt(reason);
}
\end{ccode}

\noindent \textbf{What the handler decides is what the error costs:}

\begin{itemize}
  \item \textbf{If it returns}, the kernel aborts only the thread that caused the error, and every other
    thread keeps running. A thread waiting in \code{k_thread_join()} for the aborted one returns,
    because an aborted thread counts as finished.
  \item \textbf{If it calls \code{k_fatal_halt()}}, the system halts, as with the default handler.
\end{itemize}

\textbf{Returning from a panic is not allowed.} With \code{CONFIG_ASSERT=y}, the kernel checks for
it, and the check is itself an assertion, so it panics again, and calls the same handler, which
returns again. On \code{native_sim} that recursion ends when the host process runs out of stack,
with a segmentation fault and exit status 139:

\begin{console}
ASSERTION FAIL [reason != K_ERR_KERNEL_PANIC] @ WEST_TOPDIR/deps/zephyr/kernel/fatal.c:135
	Attempted to recover from a kernel panic condition
\end{console}

\noindent \code{LOG_PANIC()} switches logging to synchronous. On a board, logging is deferred by
default (Chapter~\ref{c1:ch}): the messages wait for a thread of their own, which may never run
again after a fatal error. On \code{native_sim} logging is immediate already, so the call changes
nothing you can see there; it belongs in the handler anyway, because the handler is the same code on
a board.

A real product’s handler usually does more: it saves the reason and the thread’s name somewhere that
survives a reset, such as a retained RAM area or a setting, and resets the device, so that the next
boot can report what happened. The \file{lectures/L06/examples/debugging} example’s handler is
\file{lectures/L06/examples/debugging/src/fatal.c}.

\subsection{Assertions}
An \textbf{assertion} states something the code assumes, and checks it at run time:

\begin{ccode}
__ASSERT((percent >= 0) && (percent <= 100), "duty cycle %d %% is out of range", percent);
\end{ccode}

\noindent \code{__ASSERT_NO_MSG(condition)} is the same without a message. A failed assertion prints
the condition, where it is, and the message, and then panics:

\begin{console}
ASSERTION FAIL [(percent >= 0) && (percent <= 100)] @ WEST_TOPDIR/lectures/L06/exercises/exercise3/src/main.c:40
	duty cycle 150 % is out of range
@ WEST_TOPDIR/deps/zephyr/lib/os/assert.c:43
[00:00:00.000,000] <err> os: >>> ZEPHYR FATAL ERROR 4: Kernel panic on CPU 0
\end{console}

\noindent \textbf{Assertions are compiled out unless \code{CONFIG_ASSERT=y}}, in Zephyr’s code and
in yours. Without it the condition is not even evaluated, and the program above carries on with a
duty cycle of 150 \%. With it, CMake warns that \code{__ASSERT() statements are globally ENABLED},
because every check costs flash and time. That makes assertions a development tool: tests and debug
builds enable them, and most products ship without.

Two consequences follow:

\begin{itemize}
  \item \textbf{An assertion must not have side effects.}
    \code{__ASSERT(k_sem_take(&sem, K_NO_WAIT) == 0, "...")} takes the semaphore only in builds with
assertions.
  \item \textbf{An assertion is not error handling.} Input from a user, a sensor or the network can be wrong
    in a correct program, and has to be checked with an \code{if} in every build.
\end{itemize}

For conditions the compiler can decide, \code{BUILD_ASSERT(condition, "message")} checks at compile
time and costs nothing at run time.

\subsection{Sanitizers}
The \textbf{sanitizers} are compiler instrumentation. The compiler adds a check to every memory
access or arithmetic operation it can, and a runtime library reports every check that fails, with
the source line, the variable, and a stack trace. They find bugs that on a board corrupt memory
silently and crash, if at all, somewhere else, much later.

\subsubsection{AddressSanitizer}
\begin{shell}
west build -p -b native_sim/native/64 lectures/L06/examples/debugging \
    -- -DEXTRA_CONF_FILE=asan.conf
\end{shell}

\code{CONFIG_ASAN=y} builds with \code{-fsanitize=address}. AddressSanitizer surrounds every stack
array, heap block and global with poisoned \emph{redzones}, and checks every read and write against
them. It finds out-of-bounds reads and writes, use after free, and double frees. Exercise
Set~\ref{c6:set:4}’s battery monitor reads one element past an array on the stack:

\begin{console}
==1354852==ERROR: AddressSanitizer: stack-buffer-overflow on address 0x78a35ddf4070 ...
READ of size 4 at 0x78a35ddf4070 thread T2
    #0 0x402a6d in average .../exercise4/src/main.c:18

Address 0x78a35ddf4070 is located in stack of thread T2 at offset 48 in frame
    #0 0x402aa0 in main .../exercise4/src/main.c:33

  This frame has 1 object(s):
    [32, 48) 'window' (line 34) <== Memory access at offset 48 overflows this variable
\end{console}

\noindent Read it from the top: what kind of access (\code{READ of size 4}), where in the code (line
18, in \code{average()}), and whose memory it overflowed (\code{window}, declared on line 34).
\code{T2} is a host thread: every Zephyr thread on \code{native_sim} is one. The report ends with
\code{SUMMARY: AddressSanitizer: stack-buffer-overflow ... in average}, and the process exits with
\textbf{status 1} at the first error. (\code{CONFIG_ASAN_RECOVER}, on by default, compiles the
checks so that they \emph{could} continue, but the runtime only does so when \code{ASAN_OPTIONS}
contains \code{halt_on_error=0}.)

\textbf{AddressSanitizer exists only on \code{native_sim}}: \code{CONFIG_ASAN} depends on
\code{ARCH_POSIX}, because the checks call into the host’s \code{libasan}. The Kconfig help also
notes that leak detection at exit is off for 64-bit boards built with GCC, so
\code{native_sim/native/64} reports no leaks when the program ends.

\subsubsection{UndefinedBehaviorSanitizer}
\begin{shell}
west build -p -b native_sim/native/64 lectures/L06/examples/debugging \
    -- -DEXTRA_CONF_FILE=ubsan.conf
\end{shell}

\code{CONFIG_UBSAN=y} builds with \code{-fsanitize=undefined}, which checks for undefined behaviour
in arithmetic and pointers: signed integer overflow, shifts by too much, division by zero,
misaligned pointers, and out-of-range enum and bool values. Exercise Set~\ref{c6:set:4}’s power
calculation overflows an \code{int32_t}:

\begin{console}
.../exercise4/src/main.c:26:32: runtime error: signed integer overflow: 2250000 * 1000 cannot be represented in type 'int'
5000 mV, average 3725 mV, -2044967296 uW
\end{console}

\noindent Two things differ from AddressSanitizer:

\begin{itemize}
  \item \textbf{The program continues}, with whatever the undefined operation produced, here a wrapped
    negative number. One run can report several problems.
  \item \textbf{The report goes to stderr}, not to the console’s stdout, so a tool that only reads stdout
    never sees it. Twister’s console harness is one such tool.
\end{itemize}

UndefinedBehaviorSanitizer is not tied to \code{native_sim} as tightly. Its reports need a runtime
library, the host’s or picolibc’s; on other targets \code{CONFIG_UBSAN_TRAP} turns undefined
behaviour into a CPU exception, without saying what or where.

\textbf{Both can be enabled at once} (\code{-DCONFIG_ASAN=y -DCONFIG_UBSAN=y}), and a program that
is correct prints exactly the same with them as without. That is how Exercise Set~\ref{c6:set:4} is
tested.

\subsubsection{What a program without them does}
The same battery monitor, built without sanitizers, printed this on the machine these notes were
written on:

\begin{console}
3300 mV, average 249333220 mV, 1485000000 uW
5000 mV, average 249333645 mV, -2044967296 uW
\end{console}

\noindent The out-of-bounds read picked up whatever was next to the array, and the overflow wrapped.
Neither crashed. Another build of the same code averaged \code{261933924 mV} instead, and on another
compiler, or with another optimization level, the average could as easily have looked plausible.
That is the case for the sanitizers: undefined behaviour does not announce itself.

\subsection{GDB}
A \code{native_sim} program is an ordinary Linux executable with debug information, so GDB debugs it
directly, with no probe and no GDB server:

\begin{shell}
sudo apt install gdb
gdb --args ./build/zephyr/zephyr.exe -no-rt
\end{shell}

\begin{aside}
\textbf{Not verified here.} GDB was not installed on the machine these notes were written on, so the
commands in this section were not run while writing them. They are GDB’s standard commands, and the
Kconfig options they rely on were checked in Zephyr’s source.
\end{aside}

The commands that matter most:

\begin{booktable}{@{}LL@{}}
  \thead{Command} & \thead{Effect} \\
  \midrule
  \code{break average}, \code{break main.c:18} & Stop at a function, or at a line. \\
  \code{break k_sys_fatal_error_handler} & Stop at every fatal error, before the handler decides. \\
  \code{run}, \code{continue} & Start the program, or carry on after a stop. \\
  \code{backtrace} & The calls that led to where the program stopped. \\
  \code{print window[i]}, \code{info locals} & Values of variables. \\
  \code{info threads}, \code{thread 2} & The host threads, and switching between them. \\
\end{booktable}

Two Zephyr options make it easier:

\begin{itemize}
  \item \textbf{\code{CONFIG_DEBUG_OPTIMIZATIONS=y}} compiles with \code{-Og}, which the
    \file{lectures/L06/examples/debugging} example sets, so that variables are not optimized away and
    lines run in the order they are written.
  \item \textbf{\code{CONFIG_ARCH_POSIX_TRAP_ON_FATAL=y}} raises \code{SIGTRAP} in the halt path, before the
    process exits, so a debugger stops at a fatal error without a breakpoint.
\end{itemize}

\code{native_sim} creates a host thread for every Zephyr thread when the Zephyr thread is created,
next to the simulator’s own. A thread that is stalled in a busy loop (Chapter~\ref{c3:ch}) is easy
to find this way: interrupt the program with \code{Ctrl+C} in GDB, and look at \code{info threads}
and \code{backtrace}.

\subsection{Coverage}
\textbf{Coverage} measures which lines and branches of the code a run executed. On \code{native_sim}
it is GCC’s own gcov:

\begin{shell}
west build -p -b native_sim/native/64 lectures/L06/exercises/exercise4 \
    -- -DCONFIG_COVERAGE=y
./build/zephyr/zephyr.exe -no-rt -stop_at=1
\end{shell}

\noindent With \code{CONFIG_COVERAGE=y} the program counts every line it runs, and \textbf{writes
the counts when it exits}: \code{-stop_at} is an exit, and so is a fatal error’s. They go into a
\code{.gcda} file next to each object file, such as \code{build/CMakeFiles/app.dir/src/main.c.gcda},
which \code{gcov} reads:

\begin{shell}
cd build
gcov -o CMakeFiles/app.dir/src/main.c.gcda \
    ../lectures/L06/exercises/exercise4/src/main.c
\end{shell}

\begin{console}
Lines executed:100.00% of 14
Creating 'main.c.gcov'
\end{console}

\noindent (The number of lines depends on how your \code{exercise4} is written.)

\code{main.c.gcov} is the source with a count beside every line:

\begin{console}
       24:   18:	for (size_t i = 0; i < n; i++) {
       18:   19:		sum += window[i];
\end{console}

\noindent The loop body ran 18 times over the six calls, \code{1 + 2 + 3 + 4 + 4 + 4}, and its
condition was tested 24 times, once more per call, when it failed.

Twister collects coverage over a whole test run, and turns it into an HTML report:

\begin{shell}
sudo apt install lcov
west twister -T lectures/L06/examples/debugging -p native_sim/native/64 \
    --coverage --coverage-tool lcov
\end{shell}

\noindent The report is \code{twister-out/coverage/index.html}. Twister’s default tool is
\code{gcovr} instead, which is installed with \code{pip install gcovr}.

On a board, coverage is possible too (\code{CONFIG_COVERAGE_GCOV}), but the counts have to be dumped
over a serial line and converted on the host, and the instrumentation needs RAM that small devices
do not have. On \code{native_sim} it is one option.

\subsection{Summary}
\begin{itemize}
  \item A fatal error has a reason: an oops means a thread cannot go on, a panic that the system cannot.
    Assertions panic.
  \item The kernel logs every fatal error, and calls \code{k_sys_fatal_error_handler()}. The default one
    halts; on \code{native_sim} that is exit status 1.
  \item A handler of your own decides what an error costs. Returning aborts only the faulting thread; a
    panic must halt, with \code{k_fatal_halt()}.
  \item \code{__ASSERT()} is compiled out unless \code{CONFIG_ASSERT=y}. Assertions must have no side
    effects, and are not error handling.
  \item AddressSanitizer finds out-of-bounds access and use after free, stops at the first error, and exists
    only on \code{native_sim}.
  \item UndefinedBehaviorSanitizer finds undefined arithmetic, reports on stderr, and continues.
  \item A \code{native_sim} program is debugged with GDB like any Linux program, one host thread per Zephyr
    thread.
  \item \code{CONFIG_COVERAGE=y} writes gcov data when the program exits; Twister’s \code{--coverage} turns
    a test run’s into a report.
\end{itemize}
